\section{Related Work}

The rapid advancement of large language models (LLMs) has dramatically increased the quality and prevalence of AI-generated text, making it increasingly difficult to distinguish from human-authored content. This challenge has spurred significant research in two interconnected domains: detecting AI-generated text and identifying fake news and misinformation. Both areas require large-scale, diverse, and well-annotated datasets to develop robust detection methods. In this section, we survey existing datasets and methods in these two critical research areas, positioning our contribution within this broader landscape.

\subsection{AI-Generated Text Detection}

\subsubsection{Detection Datasets}

The development of effective AI-generated text detectors relies heavily on high-quality benchmark datasets. Several recent efforts have contributed valuable resources to this domain.

\textbf{RAID} \cite{raid2024} represents the largest and most challenging benchmark for machine-generated text detection, containing over 6 million text generations spanning 11 different language models across 8 diverse domains including arXiv abstracts, recipes, Reddit posts, news articles, poetry, and Wikipedia. The benchmark includes 11 adversarial attack types and 4 decoding strategies, providing comprehensive evaluation scenarios. Research using RAID has demonstrated that current detectors are easily fooled by adversarial manipulations, variations in sampling strategies, and unseen generative models, highlighting the need for more robust detection approaches.

\textbf{M4} \cite{m42024} is a multi-generator, multi-domain, and multi-lingual black-box dataset that includes text generated by GPT-4, ChatGPT, GPT-3.5, Cohere, Dolly-v2, and BLOOMz across seven languages. The dataset encompasses diverse domains including Wikipedia, WikiHow, Reddit, arXiv, and PeerRead, with notable coverage of low-resource languages such as Arabic, Urdu, Indonesian, and Bulgarian. This dataset won the Best Resource Paper Award at EACL 2024 and has been instrumental in demonstrating the difficulty of detector generalization across unseen domains and models.

\textbf{TuringBench} \cite{turingbench2022} provides a benchmark environment containing approximately 10,000 news articles, predominantly focused on political content, with synthetic counterparts generated by 19 different language models including various GPT-2 variants, GPT-3, GROVER models, CTRL, XLM, XLNET, and Transformer-XL. The dataset uses real news titles as prompts for generation, creating a realistic evaluation scenario for detection systems in the news domain.

\textbf{HC3 (Human ChatGPT Comparison Corpus)} \cite{hc32023} is the first large-scale human-ChatGPT comparison dataset, containing tens of thousands of comparison responses from human experts and ChatGPT across open-domain, financial, medical, legal, and psychological areas. Created shortly after ChatGPT's launch, HC3 provides direct human-AI comparison pairs that facilitate research on distinguishing ChatGPT outputs from human expert responses. An extended version, HC3 Plus, incorporates semantic-invariant tasks such as summarization, translation, and paraphrasing.

Other notable datasets include \textbf{LLM-DetectAIve} \cite{llmdetectaive2024}, which provides 236,000 examples with fine-grained labels distinguishing machine-humanized and human-polished texts across arXiv, Wikipedia, Reddit, and student essay domains; \textbf{FAIDSet} \cite{faidset2024}, an 84,000-text multilingual dataset focusing on diverse forms of human-LLM collaborative generation; and \textbf{DetectRL} \cite{detectrl2024}, a real-world benchmark accepted at NeurIPS 2024 that emphasizes ecological validity in detection scenarios.

\subsubsection{Detection Methods}

Detection methods for AI-generated text can be broadly categorized into zero-shot approaches, rewriting-based techniques, watermarking systems, and supervised classifiers.

\textbf{Zero-shot methods} require no training on labeled data. DetectGPT \cite{mitchell2023detectgptzeroshotmachinegeneratedtext} exploits probability curvature, demonstrating that AI-generated text typically lies near regions of negative curvature in the log probability space of language models. Binoculars \cite{hans2024spotting} achieves over 90\% accuracy by computing the ratio of perplexity to cross-entropy between two similar language models, showing particularly strong performance at low false positive rates in the RAID benchmark evaluation.

\textbf{Rewriting-based approaches} leverage the observation that LLMs exhibit "stubbornness" when rewriting text. RAIDAR \cite{mao2024raidargenerativeaidetection}---the foundation of our baseline method---prompts an LLM to rewrite input text and measures edit distance, exploiting the tendency of LLMs to make fewer modifications to AI-generated content than to human-written text. This approach surpasses previous methods by up to 29\% in detection accuracy. Fast-DetectGPT \cite{fastdetectgpt2024} optimizes the original DetectGPT approach, achieving 340× faster inference while maintaining comparable accuracy.

\textbf{Watermarking techniques} embed detectable signatures during text generation. Google's SynthID \cite{han2025robustness} modifies token probability distributions (logits) to embed a pseudorandom watermark that is algorithmically verifiable with a secret key, though it remains vulnerable to paraphrasing attacks. SynGuard \cite{synguard2024} extends watermarking with semantic-level guidance, achieving an 11\% improvement in F1 score against paraphrasing and translation attacks by embedding watermarks that survive meaning-preserving transformations.

\textbf{Supervised classifiers} include fine-tuned RoBERTa models \cite{roberta_openai} that achieve high in-distribution accuracy but suffer from poor out-of-distribution generalization. RADAR \cite{radar_openreview} addresses this limitation through adversarial training against paraphrasers, significantly improving robustness to meaning-preserving text transformations.

\subsection{Fake News Detection}

\subsubsection{Fake News Datasets}

Parallel to AI-generated text detection, fake news detection has produced numerous benchmark datasets addressing content authenticity and veracity.

\textbf{LIAR} \cite{liar2017} is a seminal dataset containing 12,836 manually labeled short statements collected over a decade from PolitiFact.com. Each statement is labeled with one of six fine-grained truthfulness ratings: pants-fire, false, barely-true, half-true, mostly-true, and true. With an average statement length of only 19 words, LIAR provides a challenging benchmark for claim-level fact-checking and was an order of magnitude larger than previous fake news datasets at the time of its release.

\textbf{LIAR-PLUS} \cite{liarplus2019} extends the original LIAR dataset with approximately 10,000 claims augmented with journalist-written justifications and explanations, providing valuable context for understanding fact-checking rationales in the political domain.

\textbf{FakeNewsNet} \cite{fakenewsnet2020} is a comprehensive multi-dimensional data repository containing news content, social context, and spatiotemporal information. The dataset includes two subsets: PolitiFact (434 real and 367 fake news articles) and GossipCop (for entertainment news verification). With an average article length of 1,459 words---significantly longer than LIAR---FakeNewsNet enables research on document-level fake news detection and incorporates rich social network features for studying information diffusion patterns.

\textbf{FEVER (Fact Extraction and VERification)} \cite{fever2018} contains 185,445 claims manually generated by altering sentences from Wikipedia introductory sections. Each claim is classified as Supported, Refuted, or NotEnoughInfo, with claims requiring composition of evidence from multiple sentences in 16.82\% of cases. Annotators achieved a Fleiss kappa of 0.6841, indicating substantial inter-annotator agreement. FEVER has become a foundational benchmark for evidence-based fact verification research.

\textbf{MultiFC} \cite{multifc2019} provides a multi-domain fact-checking corpus covering diverse topics beyond politics, enabling research on cross-domain generalization. \textbf{MMCFND} \cite{mmcfnd_arxiv} addresses the gap in low-resource language research with a multimodal multilingual dataset incorporating caption-aware fake news detection for Indic languages, including text-image pairs to study multimodal consistency.

\textbf{Fakeddit} \cite{fakeddit2020} is a large-scale Reddit-based dataset incorporating image modality and social media context for multimodal fake news detection. \textbf{PHEME} \cite{pheme2016} focuses on Twitter-based rumor detection and verification, providing a valuable resource for social media misinformation research.

Domain-specific datasets addressing health misinformation include \textbf{CoAID} \cite{coaid2020} and \textbf{CONSTRAINT} \cite{constraint2021}, both targeting COVID-19 misinformation in social media posts, enabling research on real-time misinformation tracking during public health crises.

\subsubsection{Detection Methods}

Fake news detection methods employ diverse approaches including fact-checking, multimodal analysis, social context integration, and knowledge graph-based reasoning.

\textbf{Fact-checking approaches} focus on evidence retrieval and claim verification, typically decomposing the task into evidence extraction from knowledge sources followed by entailment or contradiction detection. \textbf{Multimodal analysis} methods examine text-image consistency to identify manipulated or misleading content where visual and textual information contradict each other. \textbf{Social context integration} approaches analyze network propagation patterns, user credibility signals, and engagement dynamics to identify suspicious information diffusion. \textbf{Knowledge graph-based methods} leverage external structured knowledge to enhance reasoning capabilities and ground detection systems in verifiable facts.

\subsection{Positioning Our Dataset}

While existing datasets have advanced both AI-generated text detection and fake news research, significant gaps remain. Most AI-generated text detection datasets rely on synthetic prompts, student essays, or generic web content rather than high-quality journalistic sources. Fake news datasets, though valuable, often focus on short claims or social media posts rather than full-length articles. Furthermore, few datasets systematically compare outputs from multiple state-of-the-art LLMs under controlled conditions.

Our dataset addresses these limitations through six key differentiators: (1) \textbf{Source quality}---we utilize authentic New York Times articles spanning over two decades (2000-present), providing real-world, professionally authored journalistic content rather than synthetic prompts or student writing; (2) \textbf{Temporal span}---the 20+ year coverage captures evolving writing styles, topics, and news landscapes; (3) \textbf{Model diversity}---we systematically generate content using six state-of-the-art LLMs (Gemma-2-9b, Mistral-7B, Qwen-2-72B, LLaMA-8B, Yi-Large, and GPT-4-o), enabling comprehensive cross-model analysis; (4) \textbf{Scale and balance}---with 58,502 total samples and approximately 7,300 per source, the dataset provides balanced representation across all generators; (5) \textbf{Dual-task support}---uniquely, our dataset enables both binary classification (human vs. AI, Task A) and fine-grained model attribution (identifying the specific LLM generator, Task B); and (6) \textbf{Rich metadata}---each sample includes the original article abstract as a controlled prompt alongside the full human-authored narrative, facilitating nuanced analysis of generation quality and detectability.

By bridging real-world journalism with modern LLM-generated content, our dataset establishes a rigorous benchmark for advancing AI-generated text detection in high-stakes domains where content authenticity and journalistic integrity are paramount. The challenging baseline results we present---58.35\% accuracy for human vs. AI detection and 8.92\% for model attribution---underscore the difficulty of these tasks and the substantial room for methodological innovation in this critical research area.
